\section{Weyl--Heisenberg quantization over finite local rings}
\label{sec:local-rings}

This section records the first extension from finite fields to finite
commutative local rings.  In the preferred three-stage pipeline (P1), the
ring first produces the symplectic module $R^2$; a character and a polarizing
cocycle then produce its Weyl algebra and Heisenberg group; finally the
irreducible unitary representations are classified.  The decisive condition
is that the chosen character generate the character module.  At odd residue
characteristic this gives a single irreducible representation of dimension
$|R|$.

There is a genuine fork at the first stage.  Route A, used here, takes
$R\oplus R$ with an $R$-valued alternating form and then composes with a
chosen additive character.  Route B would take
$R\oplus\widehat R$ with its canonical circle-valued Pontryagin pairing; it
needs no character to become nondegenerate.  Route B and the comparison of
the physics retained by the two routes are deferred; no result below ranks
the routes.

\subsection{The local datum and its phase space}

\begin{definition}[Finite commutative local ring datum]
\label{def:finite-local-datum}
A finite local ring datum is $(R,\mathfrak m,\kappa,q)$ with $R$ a finite
commutative unital local ring with $1\neq0$, $\mathfrak m$ its unique maximal
ideal, $\kappa:=R/\mathfrak m$, and $q:=|\kappa|$; write $R^\times$ for its
unit group.  For an ideal $I\subseteq R$, put
\[
 \operatorname{Ann}(I):=\{r\in R:rI=0\},
 \qquad
 \operatorname{soc}(R):=\operatorname{Ann}(\mathfrak m).
\]
These are stipulations.  Nilpotence of $\mathfrak m$, nonvanishing of the
socle, and the other finite-local structure properties are established
below, not assumed here.  The definition is uniform in the residue
characteristic.
\end{definition}

\begin{scope}
Only finite commutative unital local rings with $1\neq0$ are defined here.
No Frobenius, characteristic, principal-ideal, or chain-ring hypothesis is
part of the datum.
\end{scope}

\provenance{D12}{---}{---}{theory/verdicts/fcr-local-r1.md}

\begin{definition}[Characters and generating characters]
\label{def:local-generating-character}
Let
\[
 \widehat R:=\operatorname{Hom}((R,+),\mathbb C^\times),
 \qquad X(R):=\widehat R\setminus\{1\}.
\]
For $\psi\in X(R)$ and $u\in R$, put $\psi_u(x):=\psi(ux)$.  Define
\[
 I_\psi:=\sum\{I\subseteq R:I\text{ is an ideal and }I\subseteq\ker\psi\},
 \qquad
 \operatorname{Gen}(R):=\{\psi\in X(R):I_\psi=0\}.
\]
Thus $I_\psi$ is the largest ideal contained in $\ker\psi$, and a member of
$\operatorname{Gen}(R)$ is a \emph{generating character}.  This definition
is uniform in the residue characteristic; existence is not stipulated.
\end{definition}

\begin{scope}
The definition includes every nontrivial additive character, not only a
trace-normalized family.  It makes no assertion that generating characters
exist for an arbitrary finite local ring.
\end{scope}

\provenance{D13}{---}{---}{theory/verdicts/fcr-local-r1.md}

\begin{definition}[Symplectic module and phase perpendicularity]
\label{def:local-symplectic-module}
Put $V(R):=R\oplus R$ and
\[
 \omega\big((a,b),(a',b')\big):=ab'-a'b.
\]
Given $\psi\in X(R)$, write
\[
 B_\psi(v,w):=\psi(\omega(v,w)),
 \qquad
 L^{\perp_\psi}:=\{v\in V(R):B_\psi(v,l)=1
                         \text{ for every }l\in L\}
\]
for an $R$-submodule $L\subseteq V(R)$.  A \emph{Lagrangian} is an
$R$-submodule satisfying $L=L^{\perp_\psi}$, and
$\operatorname{rad}(B_\psi):=V(R)^{\perp_\psi}$.  Bilinearity, strong
alternation, $R$-nondegeneracy of $\omega$, and nondegeneracy of $B_\psi$
are claims, not stipulations; the character-valued pairing $B_\psi$, not
$R$-nondegeneracy alone, is load-bearing.  This definition is uniform in
the residue characteristic.
\end{definition}

\begin{scope}
Perpendicularity is defined only for $R$-submodules of $R^2$ and is measured
through the chosen character.  No freeness condition is imposed on a
Lagrangian.
\end{scope}

\provenance{D14}{---}{---}{theory/verdicts/fcr-local-r1.md}

\begin{definition}[Admissible polarizing cocycles over a local ring]
\label{def:local-polarizing-cocycle}
For the symplectic module above, set
\[
\begin{split}
 \operatorname{Adm}(\omega)
   &:=\{\beta:V(R)\times V(R)\to R:
          \beta\text{ is $R$-bilinear and }\beta-\beta^T=\omega\},\\
 \operatorname{Sym}_R(V(R))
   &:=\{s:V(R)\times V(R)\to R:
          s\text{ is $R$-bilinear and }s=s^T\}.
\end{split}
\]
A \emph{polarizing cocycle} is a named choice
$\beta\in\operatorname{Adm}(\omega)$.  The reference member is
\[
 \beta_0\big((a,b),(a',b')\big):=ab'.
\]
This is the nonsymmetrized convention: no inverse of $2$ is assumed.  The
definition is uniform in the residue characteristic.
\end{definition}

\begin{scope}
The reference member is a stipulated coordinate formula, not a claim that
the torsor has a canonical point.  The symmetrized member $\omega/2$ is used
only under the explicit hypothesis $2\in R^\times$.
\end{scope}

\provenance{D15}{---}{---}{theory/verdicts/fcr-local-r1.md}

\begin{definition}[Local Weyl algebra, Heisenberg group, and models]
\label{def:local-weyl-system}
For the data above define
\[
 A_{\psi,\beta}(V(R)):=\bigoplus_{v\in V(R)}\mathbb C\,W_\beta(v),
 \qquad
 W_\beta(v)W_\beta(v')
   :=\psi(\beta(v,v'))W_\beta(v+v'),
\]
and
\[
 H_\beta(R):=R\times V(R),\qquad
 (t,v)(t',v'):=(t+t'+\beta(v,v'),v+v').
\]
For an $R$-submodule $L$, put
$A_L:=\operatorname{span}_{\mathbb C}\{W_\beta(l):l\in L\}$.  Given a unital
algebra character $\chi:A_L\to\mathbb C$, put
\[
 M_{L,\chi}:=A_{\psi,\beta}(V(R))
             \otimes_{A_L}\mathbb C_\chi.
\]
The fixed reference model is
$\ell^2(R)=\bigoplus_{y\in R}\mathbb C e_y$, with $\{e_y\}$ orthonormal,
\[
 X(a)e_y:=e_{y+a},\qquad Z(b)e_y:=\psi(by),
\]
and, exactly,
\[
 W_{\beta_0}(a,b):=Z(-b)X(a),\qquad
 W_{\beta_0}(a,b)e_y=\psi(-b(y+a))e_{y+a}.
\]
The sign and ordering are stipulations.  Every displayed formula is uniform
in the residue characteristic and uses no half.
\end{definition}

\begin{scope}
The algebra and group depend on the named cocycle, while the algebra and
representation also depend on the named character.  A submodule and its
character label a realization; neither is added to the underlying algebra
datum.
\end{scope}

\provenance{D16}{---}{theory/checks/fcr\_local\_check.py (G4)}{theory/verdicts/fcr-local-r1.md}

\subsection{Generating characters and the collapse boundary}

\statusProved
\begin{theorem}[Local Frobenius criterion and the unit torsor]
\label{thm:local-frobenius-criterion}
Let $R$ be a finite commutative local ring with maximal ideal
$\mathfrak m$ and residue field $\kappa$.  Then
\[
 \operatorname{Gen}(R)\neq\varnothing
 \quad\Longleftrightarrow\quad
 \dim_\kappa\operatorname{soc}(R)=1.
\]
If $\psi\in\operatorname{Gen}(R)$, then $\psi_u$ is generating exactly when
$u\in R^\times$, and $R^\times$ acts freely and transitively on
$\operatorname{Gen}(R)$.
\end{theorem}

\begin{scope}
No odd-characteristic hypothesis is used.  For rings whose socle has
$\kappa$-dimension greater than one, the theorem asserts nonexistence of a
generating character; it does not propose a replacement quantization.
\end{scope}

\provenance{D12, D13, FCR-GEN}{theory/fcr-local.md \S1}{theory/checks/fcr\_local\_check.py (G1,G8)}{theory/verdicts/fcr-local-r1.md}

\begin{proof}
Finiteness makes $R$ Artinian, so $\mathfrak m$ is nilpotent.  Wood's module
socle $S({}_RR)$ equals $\operatorname{Ann}(\mathfrak m)$: every simple ideal
$J$ has $\mathfrak mJ=0$, since $\mathfrak mJ=J$ would contradict
nilpotence; conversely, if $0\neq x\in\operatorname{Ann}(\mathfrak m)$, then
$Rx=\kappa x$ is one-dimensional over $\kappa$ and is simple.  Wood's local
Frobenius condition is therefore
$\kappa\cong\operatorname{soc}(R)$.

Wood's character-module theorem and kernel-ideal criterion identify this
condition with the existence of $\psi$ whose kernel contains no nonzero
ideal, namely $I_\psi=0$.  For such a $\psi$, multiplication by a unit
preserves $I_\psi=0$, whereas a nonunit lies in $\mathfrak m$ and kills the
nonzero socle.  The character-module isomorphism then shows that every
generating character is uniquely $\psi_u$ for a unit $u$.
\end{proof}

\statusProved
\begin{theorem}[Radical of the character-valued symplectic pairing]
\label{thm:local-phase-radical}
Let $R$ be a finite commutative local ring and let $\psi\in X(R)$.  The form
$\omega$ on $R^2$ is $R$-bilinear, strongly alternating, and
$R$-nondegenerate, while
\[
 \operatorname{rad}(B_\psi)=I_\psi\oplus I_\psi.
\]
Consequently $B_\psi$ is nondegenerate if and only if
$\psi\in\operatorname{Gen}(R)$.

For
$R=\mathbb F_3[x,y]/(x,y)^2$, write
$\psi_{a,b,c}(r+sx+ty)=\zeta_3^{ar+bs+ct}$.  The two nontrivial characters
with $(b,c)=(0,0)$ have $I_\psi=(x,y)$, and each of the remaining $24$ has
$I_\psi=\mathbb F_3(cx-by)$.  Thus all $26$ nontrivial characters collapse
the pairing through a proper quotient.
\end{theorem}

\begin{scope}
The radical formula is uniform in the residue characteristic.  The
square-zero-ring census records the obstruction at one designated seed only;
it does not construct a quantum system for a degenerate pairing and does not
state a general quotient-collapse theorem.
\end{scope}

\provenance{D12, D13, D14, FCR-RAD}{theory/fcr-local.md \S2}{theory/checks/fcr\_local\_check.py (G2,G3); exhaustive 26-character seed census in G1}{theory/verdicts/fcr-local-r1.md}

\begin{proof}
Testing $\omega((a,b),-)$ against $(0,r)$ and $(r,0)$ shows that a radical
vector must have $aR,bR\subseteq\ker\psi$, hence $a,b\in I_\psi$.  The reverse
inclusion follows because $I_\psi$ is an ideal.  Bilinearity and strong
alternation are immediate from the displayed determinant formula, and tests
against $(0,1)$ and $(1,0)$ prove $R$-nondegeneracy.  In the square-zero
example every subspace of $(x,y)$ is an ideal; restricting the character to
that two-dimensional space gives exactly the displayed kernel line, or all
of $(x,y)$ when the restriction vanishes.
\end{proof}

\subsection{Odd residue characteristic: algebra and representation}

\statusProved
\begin{theorem}[Weyl commutation over a finite local ring]
\label{thm:local-weyl-commutation}
Let $R$ be a finite commutative local ring, let $\psi\in X(R)$, and let
$\beta\in\operatorname{Adm}(\omega)$.  For every $v,v'\in V(R)$,
\[
 W_\beta(v)W_\beta(v')
   =\psi(\omega(v,v'))W_\beta(v')W_\beta(v),
\]
and
\[
 W_\beta(u)W_\beta(v)W_\beta(u)^{-1}
   =\psi(\omega(u,v))W_\beta(v).
\]
\end{theorem}

\begin{scope}
The identities are uniform in the residue characteristic and do not require
$\psi$ to be generating.  They do not by themselves imply simplicity of the
Weyl algebra.
\end{scope}

\provenance{D14, D15, D16, FCR-COMM}{theory/fcr-local.md \S3}{theory/checks/fcr\_local\_check.py (G2,G4)}{theory/verdicts/fcr-local-r1.md}

\begin{proof}
The two products have the same basis term $W_\beta(v+v')$.  Their scalar
ratio is
$\psi(\beta(v,v')-\beta(v',v))=\psi(\omega(v,v'))$; conjugation gives the
second identity.
\end{proof}

\statusProvedConditional
\begin{theorem}[Odd-characteristic transport of polarizing cocycles]
\label{thm:local-beta-odd}
Let $R$ be a finite commutative local ring with $2\in R^\times$.  Then
$\operatorname{Adm}(\omega)$ is a torsor under
$\operatorname{Sym}_R(V(R))$ and has $|R|^3$ elements.  Its unique
antisymmetric member is $\omega/2$.  For every
$\beta\in\operatorname{Adm}(\omega)$ and every symmetric $s$, the map
\[
 \varphi_s(t,v):=\left(t+\frac{s(v,v)}2,v\right)
\]
is an isomorphism $H_\beta(R)\to H_{\beta+s}(R)$ fixing the centre pointwise.
\end{theorem}

\begin{scope}
The condition $2\in R^\times$ is essential to the displayed representative
and transport.  No symmetry group of the ring-side phase space is analyzed
here.
\end{scope}

\provenance{D14, D15, D16, FCR-BETA-ODD}{theory/fcr-local.md \S4}{theory/checks/fcr\_local\_check.py (G9)}{theory/verdicts/fcr-local-r1.md}

\begin{proof}
The difference of two admissible forms is symmetric, and adding a symmetric
form preserves admissibility.  A symmetric form on $R^2$ is freely determined
by three entries.  Antisymmetry and admissibility give $2\beta=\omega$, hence
the unique member $\omega/2$.  Finally, for
$f(v)=s(v,v)/2$, symmetry gives
$f(v+w)=f(v)+f(w)+s(v,w)$, which is exactly the coboundary identity needed
for $\varphi_s$ to respect the two group laws.
\end{proof}

\statusProvedConditional
\begin{theorem}[Central simplicity and finite Stone--von Neumann theorem]
\label{thm:local-stone-von-neumann}
Let $R$ be a finite commutative local ring, let
$\psi\in\operatorname{Gen}(R)$, let
$\beta\in\operatorname{Adm}(\omega)$, and suppose $2\in R^\times$.  Then
$A_{\psi,\beta}(V(R))$ is simple, has centre $\mathbb C\,1$, has complex
dimension $|R|^2$, and is noncanonically isomorphic to
$M_{|R|}(\mathbb C)$.  The group $H_\beta(R)$ has exactly one irreducible
unitary representation with central character $\psi$, up to unitary
equivalence; its dimension is $|R|$, unitary intertwiners are unique up to
$U(1)$, and every complex algebra automorphism of the Weyl algebra is inner.
\end{theorem}

\begin{scope}
The condition $2\in R^\times$ is used to transport an arbitrary cocycle to
the fixed matrix model.  The matrix isomorphism depends on choices of
realization and basis and is not part of the algebra datum.  No
frame-preserving automorphism group is computed.
\end{scope}

\provenance{D12--D16, FCR-ALG, FCR-SVN}{theory/fcr-local.md \S5--\S6}{theory/checks/fcr\_local\_check.py (G4,G5,G6), including the G6 arena mutation}{theory/verdicts/fcr-local-r1.md}

\begin{proof}
Averaging conjugation by all Weyl basis elements kills every nonzero basis
coefficient because the pairing is nondegenerate.  It follows that the
centre is scalar and that every nonzero two-sided ideal contains $1$.
The reference operators $Z(-b)X(a)$ give a nonzero representation on
$\ell^2(R)$; simplicity makes it injective and the equal dimensions make it
onto $\operatorname{End}_{\mathbb C}(\ell^2(R))$.  The preceding transport
handles arbitrary $\beta$.

Modules with central character $\psi$ are precisely modules over this full
matrix algebra.  Its column module is the unique simple module and has
dimension $|R|$.  Schur's argument makes intertwiners scalar; polar
rescaling gives the $U(1)$ statement.  Twisting the column module by an
algebra automorphism shows that the automorphism is conjugation by an
invertible operator.
\end{proof}

\statusProvedConditional
\begin{theorem}[Lagrangians and Schrödinger models over a local ring]
\label{thm:local-lagrangians}
Let $R$ be a finite commutative local ring, let
$\psi\in\operatorname{Gen}(R)$, let
$\beta\in\operatorname{Adm}(\omega)$, and suppose $2\in R^\times$.  Every
Lagrangian $L\subseteq V(R)$ has $|L|=|R|$.  The coordinate axes are
Lagrangian, and for every ideal $I\subseteq R$ so is
\[
 I\oplus\operatorname{Ann}(I).
\]
If $R$ is not a field, then
\[
 L_*:=\operatorname{soc}(R)\oplus\mathfrak m
\]
is a nonfree Lagrangian.  Every algebra character $\chi:A_L\to\mathbb C$
produces a simple induced module $M_{L,\chi}$ of dimension $|R|$, and all
these modules are isomorphic.
\end{theorem}

\begin{scope}
This is a catalogue containing a general ideal family, not a classification
of all Lagrangians over every local ring.  In particular, it does not claim
that every Lagrangian is free or belongs to the displayed ideal family.
\end{scope}

\provenance{D12--D16, FCR-POL}{theory/fcr-local.md \S7}{theory/checks/fcr\_local\_check.py (G7), full submodule census at the two order-nine seeds}{theory/verdicts/fcr-local-r1.md}

\begin{proof}
Character orthogonality in the two orders gives
$|L|\,|L^{\perp_\psi}|=|R|^2$, so a self-perpendicular module has order
$|R|$.  Annihilator duality gives
$\operatorname{Ann}(\operatorname{Ann}(I))=I$, hence the ideal family.
For a nonfield local ring,
$\operatorname{Ann}(\mathfrak m)=\operatorname{soc}(R)$ and
$\operatorname{Ann}(\operatorname{soc}(R))=\mathfrak m$, so $L_*$ belongs
to that family.  Its quotient by $\mathfrak mL_*$ needs at least two
generators, whereas a free module of order $|R|$ would have rank one.

The restriction of $\beta$ to a Lagrangian is symmetric, and
$\chi_0(W(l))=\psi(-\beta(l,l)/2)$ is an algebra character.  The Weyl basis
makes the algebra right-free over $A_L$ on representatives for $V(R)/L$, so
the induced module has dimension $|R|$; uniqueness of the simple module
then proves that all labels give isomorphic realizations.
\end{proof}

\statusProvedConditional
\begin{proposition}[The two presentation data and bare-algebra ambiguity]
\label{prop:local-choice}
Let $R$ be a finite commutative local ring with $2\in R^\times$ and
$\operatorname{Gen}(R)\neq\varnothing$.  Beyond $R$, the presentation uses
exactly a generating character
$\psi\in\operatorname{Gen}(R)$ and a cocycle
$\beta\in\operatorname{Adm}(\omega)$.  Distinct generating characters give
inequivalent representations of the same group $H_\beta(R)$ because their
central characters differ.  All resulting bare algebras are isomorphic,
but the torsor of bare-algebra isomorphisms has no point fixed by every
target automorphism.
\end{proposition}

\begin{scope}
The last assertion concerns all bare-algebra isomorphisms.  It neither
computes nor identifies a frame-preserving isomorphism torsor and makes no
claim about a ring-side Weil or metaplectic symmetry.
\end{scope}

\provenance{D12--D16, FCR-CHOICE}{theory/fcr-local.md \S8}{theory/checks/fcr\_local\_check.py (G8,G9)}{theory/verdicts/fcr-local-r1.md}

\subsection{Clause-level return to an odd finite field}

\statusProvedConditional
\begin{theorem}[Clause-level compatibility with the odd-field system]
\label{thm:local-field-regression}
Let $R=\kappa$ be a finite field of odd characteristic and put
$q:=|\kappa|$.  Specialization recovers on the nose:
\begin{enumerate}[label=(\roman*)]
\item the nontrivial-character set and its $\kappa^\times$-torsor of order
      $q-1$;
\item bilinearity, strong alternation, and nondegeneracy of $\omega$, together
      with the zero radical of $\psi\circ\omega$;
\item the Weyl commutator and conjugation identities;
\item simplicity, scalar centre, dimension $q^2$, and a noncanonical
      $M_q(\mathbb C)$ realization;
\item the unique irreducible unitary representation of dimension $q$, its
      $U(1)$-unique intertwiners, and innerness of algebra automorphisms;
\item exactly $q+1$ Lagrangian lines, $q$ algebra characters per line,
      $q(q+1)$ model labels, and one module-isomorphism class;
\item the torsor of admissible cocycles of order $q^3$, its unique
      antisymmetric member $\omega/2$, and the centre-fixing transports;
\item the two named presentation data, the distinct-central-character
      obstruction, and isomorphism of the resulting bare algebras.
\end{enumerate}
The specialized definitions give the same $V(\kappa)$, form, admissibility
equation, Weyl and group laws, and fixed operator
$W_{\beta_0}=Z(-b)X(a)$ with the same sign.
\end{theorem}

\begin{scope}
This is deliberately a clause-level comparison.  It does not compare the
field theory's frame-preserving isomorphism torsor of order
$q^2|SL_2(\kappa)|$; the ring-side frame-preserving analysis is future work.
It also does not compare the field theorem's cocycle or automorphism-group
clauses, the $SL_2(\kappa)$-invariance of $\omega/2$, or the particular
polarization-dependent matrix construction used in the field proof.
\end{scope}

\provenance{D12--D16, FCR-GEN, FCR-RAD, FCR-COMM, FCR-BETA-ODD, FCR-ALG, FCR-SVN, FCR-POL, FCR-CHOICE, FCR-REG}{theory/fcr-local.md \S9}{theory/checks/fcr\_local\_check.py (G10), exact counts, form and radical tables, fixed matrices, rank, commutant, lines, half-census, and group profile}{theory/verdicts/fcr-local-r1.md}

\begin{proof}
For a field, $\mathfrak m=0$, $\operatorname{soc}(R)=R$, and the only ideals
are $0$ and $R$, so every nontrivial character is generating.  Substitution
in the preceding results proves every listed clause.  A self-perpendicular
subspace of the nondegenerate two-dimensional space is a line; nonzero
vectors modulo $\kappa^\times$ give $(q^2-1)/(q-1)=q+1$ lines, and each line
has $q$ additive characters.  No excluded symmetry or frame clause follows
from these calculations.
\end{proof}
